% =====================================================================
%  OWNER: F (YANG Xueli).  Content: LO4 -- alternative architectures and
%  the justification of the chosen one on security, cost and reliability.
%
%  Insert in report/main.tex immediately after
%      % =====================================================================
%  OWNER: F. Only F edits this file; others use Overleaf comments.
%  Content: DynamoDB schema, Streams, the notify Lambda and SNS email.
% =====================================================================
\subsubsection{Data Schema and Notification}\label{sec:data}

% TODO(F): write this part.

%  (it then becomes Section 3.4). Only A edits main.tex.
% =====================================================================
\subsection{Alternative Architectures Considered}\label{sec:alt}

A design that presents only the chosen architecture is an assertion rather than
an argument. Four options were costed and compared against the same workload ---
10{,}000 applications per month with a tenfold peak --- before VeriCloud was
selected. They are the three that a team with this brief would realistically
consider, plus the build-versus-buy alternative that any bank would ask about.

\begin{table}[htbp]
  \centering
  \scriptsize
  \caption{Architectures considered for the same workload. Cost figures for the
           two middle columns are indicative; the served columns are modelled in
           Section~\ref{sec:cost}.}
  \label{tab:alt}
  \begin{tabular}{@{}p{2.7cm}p{3.0cm}p{3.0cm}p{3.0cm}p{3.0cm}@{}}
    \toprule
     & Provisioned EC2 + RDS & Containerised (ECS) & \textbf{Serverless (chosen)} & Managed KYC service \\
    \midrule
    Cost at 10{,}000/month
      & \$232.48 & $\sim$\$130--200 & \textbf{\$15.80} & \$5{,}000--15{,}000 \\
    Cost at 100{,}000/month
      & $\sim$\$240 (flat) & $\sim$\$180--260 & $\sim$\$158 & \$50{,}000--150{,}000 \\
    Security exposure
      & guest OS and runtime patched by us & host and runtime patched by us & managed runtime, least-privilege IAM by default & biometric data leaves our boundary \\
    Failure handling
      & health checks and autoscaling written by us & orchestrator plus our retry logic & native Retry, Catch and dead-letter queue & defined by the vendor \\
    Operational burden
      & high & medium & low & lowest \\
    \bottomrule
  \end{tabular}
\end{table}

\paragraph{Why security favours serverless here.}
The security argument is about \emph{attack surface}, not about features. With
provisioned servers the team owns the guest operating system, its patch cycle and
its network exposure; every week of the project is another week of unpatched
packages. With containers the host is still ours. With the serverless design the
runtime is managed, each function executes under a role scoped to the resources
it actually touches, and the only inbound path is a 300-second presigned URL to a
private bucket --- there is no server to log into. The managed service scores
equally well on isolation but moves the biometric data outside the account, which
under a strict reading of data-minimisation is the largest possible exposure.

\paragraph{Why cost favours serverless.}
At the modelled workload the provisioned stack costs \$232.48 a month and the
serverless stack \$15.80 --- a 93.2\% reduction (Section~\ref{sec:cost}). The
difference is not a discount; it is a change in billing structure. Provisioned
capacity is paid for continuously and the workload is idle most of the time,
whereas serverless capacity is paid per invocation and costs nothing between
campaigns. This advantage is conditional: the sensitivity analysis in
Section~\ref{sec:cost:sensitivity} shows it narrowing to about 32\% at
100{,}000 applications per month and disappearing near 150{,}000, where the flat
cost of fully utilised servers finally beats per-request billing.

\paragraph{Why reliability favours serverless.}
The reliability primitive that matters here is \emph{isolation of failure}, and
it is the one that is hardest to build by hand. In the serverless design a
poisoned message is retried three times and parked in a dead-letter queue; a
failed face comparison is caught on its own branch while the OCR result survives;
an uncertain decision degrades to human review rather than to an error. On
provisioned or containerised infrastructure the same behaviour has to be written,
tested and maintained by the team --- and a team that is still learning the
platform is exactly the team least able to maintain it. Serverless also inherits
multi-availability-zone execution without configuration, whereas the provisioned
baseline only survives a zone failure if it was deliberately designed to.

\paragraph{Where the rejected options would win.}
The comparison is only credible if it states its own limits. Provisioned
infrastructure wins when a bank already owns depreciated hardware, has a
data-residency obligation that forbids the managed service, or runs sustained
high volume where reserved capacity is cheaper than per-request billing.
Containers win when a stage needs a custom runtime, a GPU, or a long-running job
that exceeds the Lambda execution limit. A managed KYC service wins when there
is no engineering team to build and operate the pipeline, and when a per-check
fee of a few dollars is small relative to the engineering it replaces --- a cost
that this model, which prices infrastructure only, does not capture. VeriCloud
was selected because this brief is the opposite of all three: a student team with
no existing infrastructure, a bursty low-to-moderate volume, and a requirement to
demonstrate the engineering rather than procure it.
